\section{Phân tích và thiết kế hệ thống}

\subsection{Kiến trúc tổng thể Vi dịch vụ}
Hệ thống gồm 5 tầng kiến trúc:
\begin{enumerate}
    \item \textbf{Tầng Ngoại vi:} Khách hàng chat qua CRM (Zalo/FB).
    \item \textbf{Tầng Cổng giao tiếp (Gateway):} API Gateway nhận Webhook.
    \item \textbf{Tầng Hàng đợi Bất đồng bộ:} RabbitMQ làm bộ đệm giảm chấn (Buffer).
    \item \textbf{Tầng Xử lý Giao dịch:} Các Async Worker thực hiện ACID Locking trên PostgreSQL.
    \item \textbf{Tầng AI Đa tác tử:} LangGraph quản lý StateGraph. Tích hợp Qdrant cho Hybrid Search.
\end{enumerate}

\subsection{Đồ thị Trạng thái Đa tác tử (State Graph)}
Mỗi phiên chat duy trì một \textbf{Shared State} (Biến trạng thái dùng chung).
Các Nodes là các Agent: \texttt{Orchestrator}, \texttt{Bus\_Agent}, \texttt{Hotel\_Agent}.
Các Edges là luồng điều khiển. Orchestrator đánh giá State và định tuyến (Conditional Edges) đến tác tử phù hợp.

\subsection{Sơ đồ tuần tự: Xử lý Đơn hàng \& ACID}
Gateway $\rightarrow$ Publish RabbitMQ $\rightarrow$ Worker Pull $\rightarrow$ \texttt{BEGIN} $\rightarrow$ \texttt{SELECT FOR UPDATE} $\rightarrow$ \texttt{UPDATE} $\rightarrow$ \texttt{COMMIT} $\rightarrow$ \texttt{ACK}.

\subsection{Sơ đồ tuần tự: Multi-Agent \& Circuit Breaker}
User $\rightarrow$ Orchestrator $\rightarrow$ Bus Agent (Tool Call DB) $\rightarrow$ Orchestrator $\rightarrow$ Hotel Agent.
Nếu LLM bị lỗi Rate Limit $\rightarrow$ Circuit Breaker cắt mạch $\rightarrow$ Fallback LLM Server 2.
